\scanpage{15}
Détermination du groupe des caractères rationnels sur $\mathbb Q$

\[
T_{\mathbb C}=D_{\mathbb C}(M),\quad M\text{ où \unsure{on a} }\pi_1\text{ opère}\quad(\pi_1=\pi_1(\mathbb Q,\mathbb C))
\]
\inserted{via un groupe quotient fini $\pi$ opérant fidèlement}

$\mathbb Z(M)$ anneau des caractères \inserted{virtuels} de $T_{\mathbb C}$, $\pi$ y opère
\[
K(T)=\mathbb Z(M)^\pi\qquad\text{Anneau de type fini sur }\mathbb Z.
\]

[N.B. $\mathbb Z(M)$ est l'anneau affine de $S=\deleted{\unsure{\mathbb Z}}D_{\mathbb Z}(M)$, sur lequel $\pi$ opère pr automorphismes de groupe. $K(T)$ est l'anneau affine de $D_{\mathbb Z}(M)/\pi$]. Les caractères \uline{effectifs} sont ceux qui prennent des valeurs $\geq0$ sur les points de $S(\mathbb R)^\circ$ ?? Comment définir une structure d'effectivité \uline{intrinsèque} ? Le procédé \unsure{standard} est de se fixer un $i_1:\mathbb G_{m\mathbb C}\mathbin{-}T_{\mathbb C}$ i.e. un $i_1^*:M\longrightarrow\mathbb Z$, avec les propriétés habituelles, et donnant donc $\mathbb Z(M)\mathbin{-}\mathbb Z(\mathbb Z)$ d'où
\diagram{\begin{tikzcd}[ampersand replacement=\&]
\mathbb G_{m\mathbb Z}\arrow[r,"\alpha"] \& S
\end{tikzcd}}
(\unsure{pr aug. aux $\pi$}\ldots)

\margin{$\lambda_1^{n_1}\lambda_2^{n_2}\mathbin{-}\lambda_r^{n_r}$}

La condition est, \deleted{\unsure{pr|pour}} pr un $\chi\in K(T)\simeq\mathbb Z(S)^\pi$, que $\chi\circ\alpha\unsure{\simeq|\in}$ \unsure{$\mathbb Z$} anneau affine de $\mathbb G_{m\mathbb Z}$ soit une fonction \uline{polynômiale} i.e. restant finie sur $\mathbb G_{m\mathbb Z}(\mathbb R)=\mathbb R^*$ pr $t\to0$.

Considérons les caractères de modules effectifs de poids donné $\beta$, \unsure{prouvons} qu'il n'y en a qu'un \unsure{nb} fini. On \inserted{\uuline{ne}} \uuline{peut} \inserted{\uuline{pas}} \unsure{simplifier|remplacer} dans cette question $\mathbb Q$ pr $\mathbb R$, \deleted{Alors} car si $T$ est de dim $\geq3$, la \unsure{notion|condition} de \unsure{partie} compacte est \uline{Fausse}.

\mspage{14}
